\section{Experiments}

This section is written as an executable evaluation contract.  The local paper
workspace contains the latest event-balanced source and unit-test
specifications, but no training-machine checkpoint, Gate O/R JSON, or
UniversalVLOG simulator rollout.  Accordingly, unavailable values remain
visible placeholders rather than being filled with results from the
incompatible legacy Qwen-OFT model.

\begin{table}[t]
\centering
\scriptsize
\setlength{\tabcolsep}{3pt}
\begin{tabular}{p{0.25\columnwidth}p{0.22\columnwidth}p{0.45\columnwidth}}
\toprule
Item & Status & Evidence boundary \\
\midrule
Universal source & Audited & Commit \texttt{4111898}; Qwen/shared DiT, native adapters, event-balanced VQ/router, option conditioning, controller. \\
CPU tests & Specified & Local runtime lacks \texttt{pytest}; no new raw test log generated in this paper pass. \\
Official checkpoint & Pending & Strict-load, full parity, GPU memory/time not present locally. \\
Gate O/R & Pending & Evaluator exists; no training-machine JSON recovered. \\
Simulator results & Pending & No UniversalVLOG RoboCasa/LIBERO rollout artifact recovered. \\
Legacy graph/critic & Excluded & Different architecture; cannot populate current results. \\
\bottomrule
\end{tabular}
\caption{Evidence status for this draft.  Source presence is not promoted to
benchmark evidence.}
\label{tab:evaluation-status}
\end{table}


\subsection{Research Questions}

We organize evaluation around five questions.
\begin{description}
    \item[RQ1: Base validity.] Does the universal wrapper strictly load the
    official Qwen/GR00T core and preserve fusion-off velocity/action outputs?
    \item[RQ2: Option causality.] Does the future-action oracle outperform
    wrong and shuffled codes under paired flow time and noise without
    degrading the fusion-off base?
    \item[RQ3: Router recovery.] Can a hard current-context router recover the
    oracle's code use and action loss when posterior information is removed?
    \item[RQ4: Temporal utility.] Does explicit persistence improve task success
    or temporal failure modes over fusion-off, per-query, and fixed-duration
    controls?
    \item[RQ5: Shared embodiment interface.] Can one Qwen/DiT/codebook support
    RoboCasa and LIBERO through native adapters without unacceptable forgetting
    or systems overhead?
\end{description}

\begin{table*}[t]
\centering
\scriptsize
\setlength{\tabcolsep}{3pt}
\begin{tabular}{p{0.05\textwidth}p{0.20\textwidth}p{0.22\textwidth}p{0.17\textwidth}p{0.28\textwidth}}
\toprule
RQ & Required comparison & Primary evidence & Destination & Falsifying outcome \\
\midrule
1 & Original base vs.\ universal fusion-off & Strict load, paired velocity/action parity, base rollout & Appendix & Missing core keys, parity error, or failed base floor \\
2 & Base/correct/wrong/shuffle & Paired FM gaps, action/velocity delta, active/effective code use, event alignment & Fig.~\ref{fig:gate-diagnostics} & Correct not better than wrong/shuffle, worse than base, or collapsed use \\
3 & Oracle vs.\ hard router & Accuracy, active/effective codes, router-oracle/base FM gaps & Fig.~\ref{fig:gate-diagnostics} & Router collapse or unacceptable action-loss gap \\
4 & Fusion-off/per-query/fixed/persistent & Paired success, duration, switches, failures & Table~\ref{tab:main-results} & Persistence matched or beaten by simpler controls \\
5 & Single vs.\ shared embodiments & Per-embodiment success/Gates, forgetting, latency & Appendix & Base forgetting or codebook dominated by embodiment \\
\bottomrule
\end{tabular}
\caption{Research questions, required evidence, and predeclared falsification.
FM denotes flow matching.}
\label{tab:rq-evidence}
\end{table*}


\subsection{Benchmarks and Protocol}

\paragraph{RoboCasa-GR1.}
The base route uses the native \(58\)-D state, \(29\)-D action, and horizon
\(16\) contract declared by the official GR00T-style checkpoint.  U0 first
reproduces a locked subset of RoboCasa-GR1 tasks before any option training.
The final task manifest, simulator revision, camera views, episode cap,
normalization key, action execution rate, and rollouts per task are
\textbf{[PROTOCOL PLACEHOLDER: ROBOCASA-MANIFEST]}.

\paragraph{LIBERO.}
The additional embodiment uses native \(7\)-D state and action adapters and is
evaluated on LIBERO-Spatial, LIBERO-Object, LIBERO-Goal, and LIBERO-10.
The final dataset version, observation view, action convention, executed chunk
prefix, simulator revision, and episode seeds are
\textbf{[PROTOCOL PLACEHOLDER: LIBERO-MANIFEST]}.

\paragraph{Pairing and checkpoint selection.}
All compared policies receive the same task/initial-state seeds.  Training seeds
are the unit of model variability; episode-level pairing is used within a seed.
Checkpoints are selected by a predeclared validation rule before test rollout.
No method is selected on the reported test average.  Exact seeds, denominators,
confidence interval method, and failed-run policy are
\textbf{[PROTOCOL PLACEHOLDER: STATISTICAL-PROTOCOL]}.

\subsection{Baselines and Controlled Variants}

The minimum comparison set isolates initialization, capacity, option identity,
routing, and persistence:
\begin{itemize}
    \item \textbf{Official flow base:} original Qwen/GR00T policy on its base
    embodiment.
    \item \textbf{Universal fusion-off:} the new wrapper with VLOG disabled;
    required to match the base action-head path.
    \item \textbf{Native-adapter base:} U0-native on a non-base embodiment with
    no option path.
    \item \textbf{Option-independent adapter:} a parameter-matched residual that
    cannot vary by option, testing added capacity.
    \item \textbf{Oracle option:} training-time posterior code, used only for
    offline upper-bound diagnostics.
    \item \textbf{Per-query router:} hard routed code with controller memory
    disabled.
    \item \textbf{Fixed duration:} hard routed code retained for a fixed number
    of queries.
    \item \textbf{Persistent VLOG:} hard router plus the \(d_{\min}/d_{\max}\)
    hysteretic controller.
    \item \textbf{Wrong/shuffled code:} paired negative interventions, not
    deployable baselines.
\end{itemize}
If U3 updates the shared action expert, a matched continued-fine-tuning control
uses the same optimizer steps, data, and shared-expert learning rate.

\subsection{Base and Systems Validation}

Before training options, the evaluation records: strict checkpoint key/shape
coverage; cached action-head parity at identical state, visual features,
\(\tau\), and noise; deterministic seeded action parity; normalization
round-trip by action dimension; variable-horizon masking; per-stage trainable
parameters; and GPU forward/backward/predict memory and time.  RoboCasa
fusion-off rollout must meet a predeclared base floor before U0-native or U1.
For mixed embodiments, batches are partitioned by embodiment, evaluated through
the shared Qwen/DiT, and reduced by observation count.

\subsection{Offline Gate O and Gate R}

The released evaluator repeats paired flow interventions on a fixed audit
subset.  For each observation we compute
\[
\begin{aligned}
\Delta_{-\!+}&=\mathcal{L}_{-}-\mathcal{L}_{+},&
\Delta_{\mathrm{sh}\!+}&=\mathcal{L}_{\mathrm{shuffle}}-\mathcal{L}_{+},\\
\Delta_{0+}&=\mathcal{L}_{0}-\mathcal{L}_{+}.&&
\end{aligned}
\]
Gate O requires positive confidence intervals for wrong/shuffled gaps, a
nonnegative oracle-vs-base gap, sufficient active codes at a minimum usage
share, sufficient effective option number, bounded dominant share, nonzero
residual, and bounded residual-to-action-token ratio.  When event balancing is
enabled, the gate also records option--event normalized mutual information
against a shuffled-code baseline; event IDs remain balancing strata rather than
semantic labels.  We also report codebook cosine similarity and
option-conditioned velocity difference.

Gate R reports hard-router accuracy against posterior IDs, router cross-entropy,
active-code count, effective option number, dominant share, router-vs-oracle
flow gap, and router-vs-base degradation.  Raw ID accuracy is interpreted
together with action loss because different codes may be behaviorally
equivalent.  All Gate statistics are reported separately by embodiment as well
as globally.

\begin{figure*}[t]
\centering
\fbox{\parbox{0.95\textwidth}{\centering
\textbf{FIGURE PLACEHOLDER F5: Gate O/R quantitative diagnostics}\\[0.45em]
\textbf{Panel (a):} paired flow loss for base, correct, wrong, shuffled, and
hard router, with 95\% confidence intervals.\\
\textbf{Panel (b):} correct-vs-wrong and correct-vs-shuffle relative gaps by
embodiment. \quad
\textbf{Panel (c):} oracle/router code usage, dominant share, and dead codes.\\
\textbf{Panel (d):} router accuracy and router-oracle flow-loss gap. \quad
\textbf{Panel (e):} fusion residual ratio and action/velocity intervention
magnitude.\\[0.25em]
All panels must be generated from
\texttt{evaluate\_universal\_gates.py} JSON; do not draw illustrative values.}}
\caption{\textbf{Required mechanism evidence.}  Gate O asks whether option
identity causally specializes the action expert without degrading the base;
Gate R asks whether the deployment router recovers that oracle behavior.
\textbf{[DATA PENDING]}}
\label{fig:gate-diagnostics}
\end{figure*}


\subsection{Rollout Evaluation}

The primary metric is simulator success, reported as successes/episodes and
percentage for every task and training seed.  To test the temporal mechanism,
we additionally record option changes per episode, decision-duration and
environment-step duration, \(d_{\max}\) reselection rate, router hysteresis
switches, code-use entropy, and episode reset correctness.  Task benefit is
credited to persistence only if persistent VLOG improves beyond fusion-off,
option-independent, per-query, and fixed-duration controls without option
collapse.

\begin{table*}[t]
\centering
\small
\setlength{\tabcolsep}{4pt}
\begin{tabular}{lcccccc}
\toprule
Method & RoboCasa & Spatial & Object & Goal & LIBERO-10 & Macro \\
\midrule
Official flow base & P & -- & -- & -- & -- & P \\
Universal fusion-off / U0-native & P & P & P & P & P & P \\
Option-independent adapter & P & P & P & P & P & P \\
Per-query routed option & P & P & P & P & P & P \\
Fixed-duration option & P & P & P & P & P & P \\
Persistent VLOG-VLA & P & P & P & P & P & P \\
\bottomrule
\end{tabular}
\caption{Simulator success under protocol-matched evaluation.  P is a result
placeholder; each cell requires successes/episodes, training seeds, uncertainty,
checkpoint rule, and raw per-episode evidence.  Dashes are not applicable.}
\label{tab:main-results}
\end{table*}


\subsection{Ablations and Cross-Embodiment Analysis}

The ablation matrix varies option token, FiLM, option-only conditioning,
counterfactual and improve losses, router persistence, \(K\), and whether the
shared DiT is jointly tuned.  Cross-embodiment analysis reports global and
per-embodiment code usage, mutual information between code and embodiment,
RoboCasa retention after LIBERO adaptation, and separate Gate O/R results.
A shared codebook is supported only if its effect is not explained entirely by
embodiment identity.

\begin{table*}[t]
\centering
\scriptsize
\setlength{\tabcolsep}{3.5pt}
\begin{tabular}{p{0.22\textwidth}cccccccc}
\toprule
Variant & VQ & Option token & FiLM & Opt.-only & CF/Imp. & Persist. & Success & Gate O/R \\
\midrule
Fusion-off base & No & No & No & -- & No & No & P & -- \\
Option-independent adapter & No & Yes & Yes & No & No & No & P & -- \\
Oracle conditioner & Yes & Yes & Yes & Yes & Yes & No & offline & P/-- \\
No option token & Yes & No & Yes & Yes & Yes & Rule & P & P/P \\
No FiLM & Yes & Yes & No & Yes & Yes & Rule & P & P/P \\
Legacy context conditioner & Yes & Yes & Yes & No & Yes & Rule & P & P/P \\
No CF / no improve & Yes & Yes & Yes & Yes & No & Rule & P & P/P \\
Per-query router & Yes & Yes & Yes & Yes & Yes & No & P & P/P \\
Persistent VLOG-VLA & Yes & Yes & Yes & Yes & Yes & Rule & P & P/P \\
\bottomrule
\end{tabular}
\caption{Required component attribution.  P denotes a verified-result
placeholder; CF/Imp.\ denotes paired counterfactual and improve losses.
The oracle row is an offline diagnostic rather than a deployable policy.}
\label{tab:ablation-results}
\end{table*}


\subsection{Qualitative and Failure Evidence}

For predeclared task/seed strata, synchronized video frames are aligned with
query ID, active code, router candidate, decision age, switch reason, and action
chunk.  We include both a supported recovery and a hierarchy-specific failure
when they exist; ambiguous failures remain unattributed.  Code colors are
identifiers, not semantic names.

\begin{figure*}[t]
\centering
\fbox{\parbox{0.95\textwidth}{\centering
\textbf{FIGURE PLACEHOLDER F6: Paired rollout and option timeline}\\[0.45em]
\textbf{Top:} matched frame strips for fusion-off base, per-query latent, and
persistent VLOG on the same task/initial state.\\
\textbf{Middle:} active option ID, router candidate, decision age, switch
reason, and action-chunk boundaries.\\
\textbf{Bottom:} paired task outcome plus action difference; include one
supported recovery and one hierarchy-specific failure, or mark the failure
unattributed.\\[0.25em]
Frames must link to video IDs and per-query JSONL records.}}
\caption{\textbf{Required rollout evidence.}  The asset must align controller
state with observable behavior and paired outcomes; option colors alone must
not be assigned semantic skill names. \textbf{[DATA PENDING]}}
\label{fig:rollout-evidence}
\end{figure*}


\subsection{Efficiency, Statistics, and Artifacts}

Matched profiling reports total/trainable/added parameters, inference latency,
peak memory, and throughput for fusion-off and persistent variants on identical
hardware, precision, image inputs, batch convention, and warm-up.  Success
comparisons use paired episode outcomes within each model seed and confidence
intervals over the predeclared aggregation unit.  Every populated cell must
trace to source commit and diff, resolved configuration, checkpoint, dataset
manifest, raw per-episode record, aggregation command, and file hash.

\textbf{[RESULT PLACEHOLDER: populate Tables~\ref{tab:main-results} and
\ref{tab:ablation-results}, Figure~\ref{fig:gate-diagnostics}, and
Figure~\ref{fig:rollout-evidence} only from verified UniversalVLOG artifacts.]}
